\SourceLeaf{242}%
\Needspace{6\baselineskip}\begin{center}\textsc{Gauss} an \textsc{Gerling}. Göttingen, 4. Februar 1844.\end{center}

\mbox{.\ .\ .\ .\ .\ .\ .\ .} So ist es zugegangen, dass ich vergessen habe, einen \textls{andern} Punkt Ihres Briefes zu beantworten, an den ich eben heute durch eine ganz zufällige Veranlassung wieder erinnert bin.

Der Titel des Buchs meines alten Universitätsfreundes \textsc{Varkas} (i. e. Wolfgang) von \textsc{Bolyai}, Vater, ist wörtlich folgender:

Tentamen juventutem studiosam in elementa matheseos purae, elementaris ac sublimioris, methodo intuitiva, evidenti\"aque huic propria, introducendi cum Appendice triplici. Auctore Professore Matheseos et Physices Chemiaeque Publ. Ordinario. Tomus primus. Mar[o]s Vásárhelyini 1832. Typis Collegii Reformatorum per Josephum, et Simeonem Kali de felsö Vist. LXVI u. 502 Seiten 8\textsuperscript{vo}. Tomus secundus 1833. XVI u. 400 S., viele Kupfertafeln. Den Namen des Verfassers finde ich nicht angegeben.

Die eine, beim ersten Bande befindliche Appendix hat den Titel:

Appendix scientiam spatii absolute veram exhibens: a veritate aut falsitate Axiomatis XI Euclidei (a priori haud unquam decidenda) independentem; adjecta ad casum falsitatis, quadratura circuli geometrica. Auctore \textsc{Johanne} \textsc{Bolyai} de eadem\HistoricalNote{*)}{*)}{\emph{of that ilk}, wie es gewöhnlich bei \textsc{Walter Scott} heisst. — [\emph{Vergl. auch S.} 225.]}, Geometrarum in exercitu Caesareo Regio Austriaco Castrensium Capitaneo.

Diese Appendix ist besonders paginirt (p. 1---26); da aber auf der Titelseite weder Ort noch Jahr angegeben ist, so präsumire ich, dass diese Schrift, \textls{welche gerade die in Rede stehende ist}, nicht für sich, sondern nur als Anhang des Werkes des Vaters ins Publicum gekommen ist.

Übrigens hat in den letzten Decennien ein Russe (\textsc{Lobatschewsky}, Staatsrath und Professor in Kasan) einen ähnlichen Weg eingeschlagen. Er nennt die Nicht-Euklidische die imaginäre Geometrie (wie Ihr ehemaliger Kollege [\textsc{Schweikart}] Astralgeometrie) und hat darüber in \textls{russischer} Sprache viele sehr ausgedehnte Abhandlungen gegeben (meistens in den \gausscyr{Записки Казанскаго Университета}, Memoiren der Kasanschen Universität), zum Theil auch in besondern Brochuren, die ich, glaube ich, alle besitze, aber ihre genaue Lecture noch verschoben habe,
